\documentclass{ctexart}
\usepackage{amsmath, amssymb}
\usepackage{geometry}
\geometry{a4paper, margin=2cm}

\title{\textbf{第3章 功和能} --- 例题解答}
\author{}
\date{}

\begin{document}
\maketitle

\begin{enumerate}

\item 已知 $m = 2\,\text{kg}$，在 $F = 12t$ 作用下由静止做直线运动。求 $t = 0 \to 2\,\text{s}$ 内 $F$ 作的功及 $t = 2\,\text{s}$ 时的功率。

\textbf{解：}
\[
a = \frac{F}{m} = 6t = \frac{\mathrm{d}v}{\mathrm{d}t} \;\Longrightarrow\; v = 3t^{2} = \frac{\mathrm{d}x}{\mathrm{d}t}
\]
\[
A = \int_{0}^{x} F\,\mathrm{d}x = \int_{0}^{t} 12t \cdot 3t^{2}\,\mathrm{d}t = \int_{0}^{2} 36t^{3}\,\mathrm{d}t = 36\cdot\frac{2^{4}}{4} = 144\,\text{J}
\]
\[
P = \vec{F}\cdot\vec{v} = 12t \cdot 3t^{2} = 36t^{3},\quad P(2) = 36 \times 8 = 288\,\text{W}
\]

\item 质量为 $10\,\text{kg}$ 的质点，在外力作用下做平面曲线运动，该质点的速度为 $\vec{v} = 4t^{2}\,\vec{i} + 16\,\vec{j}$，开始时质点位于坐标原点。求在质点从 $y = 16\,\text{m}$ 到 $y = 32\,\text{m}$ 的过程中，外力做的功。

\textbf{解：}
\[
v_x = \frac{\mathrm{d}x}{\mathrm{d}t} = 4t^{2},\quad v_y = \frac{\mathrm{d}y}{\mathrm{d}t} = 16
\]
由 $v_y = 16$ 得 $y = 16t$，故 $y = 16\,\text{m}$ 时 $t = 1\,\text{s}$，$y = 32\,\text{m}$ 时 $t = 2\,\text{s}$。
\[
F_x = m\frac{\mathrm{d}v_x}{\mathrm{d}t} = 10 \times 8t = 80t,\quad F_y = m\frac{\mathrm{d}v_y}{\mathrm{d}t} = 0
\]
\[
A = \int F_x\,\mathrm{d}x + F_y\,\mathrm{d}y = \int_{1}^{2} 80t \cdot (4t^{2}\,\mathrm{d}t) = \int_{1}^{2} 320t^{3}\,\mathrm{d}t = 320\cdot\frac{2^{4} - 1^{4}}{4} = 1200\,\text{J}
\]

\item 已知用力 $\vec{F}$ 缓慢拉质量为 $m$ 的小球，$\vec{F}$ 保持方向不变。求 $\theta = \theta_0$ 时 $\vec{F}$ 作的功。

\textbf{解：}缓慢拉意味着小球始终处于平衡状态。
\[
\begin{cases}
F - T\sin\theta = 0 \\
T\cos\theta - mg = 0
\end{cases}
\;\Longrightarrow\; F = mg\tan\theta
\]
元功 $\mathrm{d}A = \vec{F}\cdot\mathrm{d}\vec{r} = F\cos\theta\,\mathrm{d}s$，而 $\mathrm{d}s = L\,\mathrm{d}\theta$，故
\[
A = \int_{0}^{\theta_0} mg\tan\theta \cdot \cos\theta \cdot L\,\mathrm{d}\theta = mgL\int_{0}^{\theta_0} \sin\theta\,\mathrm{d}\theta = mgL(1 - \cos\theta_0)
\]

\item（书中例题 3.2，p.98，重点）一条长 $L$、质量 $M$ 的均匀柔绳，$A$ 端挂在天花板上，自然下垂，将 $B$ 端沿铅直方向提高到与 $A$ 端同高处。求该过程中重力所作的功。

\textbf{解：}提起高度 $y$ 时，被提起部分长度为 $\dfrac{y}{2}$，质量为 $\dfrac{M}{L}\cdot\dfrac{y}{2}$。
元功：
\[
\mathrm{d}A = -\frac12\frac{M}{L}gy\,\mathrm{d}y
\]
总功：
\[
A = \int_{0}^{L} -\frac12\frac{M}{L}gy\,\mathrm{d}y = -\frac12\frac{M}{L}g\cdot\frac{L^{2}}{2} = -\frac14 MgL
\]
负号表示重力作负功。

\item（重点）蓄水池面积 $S$，水深 $h$，水面距地面 $H$，水的密度为 $\rho$。求抽出水需要作多少功？

\textbf{解：}离地面 $x$ 处，深 $\mathrm{d}x$ 的一层水质量 $\mathrm{d}m = \rho S\,\mathrm{d}x$，将其提到地面需克服重力做功：
\[
\mathrm{d}A = gx\,\mathrm{d}m = \rho S g x\,\mathrm{d}x
\]
总功：
\[
A = \int_{H}^{H+h} \rho S g x\,\mathrm{d}x = \rho S g\cdot\frac12 x^{2}\Big|_{H}^{H+h}
= \rho S g\cdot\frac12[(H+h)^{2} - H^{2}]
= \rho g S H h + \frac12\rho g S h^{2}
\]

\item 风力 $F$ 作用于向北运动的船，风力方向变化的规律是 $\theta = Bs$，其中 $s$ 为位移，$B$ 为常数，$\theta$ 为 $F$ 与 $S$ 间的夹角。如果运动中，风的方向自南变到东，求风力作的功。

\textbf{解：}元功 $\mathrm{d}A = F\cos\theta\,\mathrm{d}s$。由 $\theta = Bs$，$\mathrm{d}s = \dfrac{\mathrm{d}\theta}{B}$。
风向自南变到东：$\theta$ 从 $0$ 变到 $\dfrac{\pi}{2}$。
\[
A = \int_{0}^{\pi/2} \frac{F\cos\theta}{B}\,\mathrm{d}\theta = \frac{F}{B}\sin\theta\Big|_{0}^{\pi/2} = \frac{F}{B}
\]

\item（书中例题 3.5，p.103）物体的质量为 $m$，弹簧劲度系数为 $k$，$A$ 板及弹簧质量不计。求自弹簧原长 $O$ 处，突然无初速度地加上物体时，弹簧的最大压缩量。

\textbf{解：}设最大压缩量为 $\lambda_{\text{max}}$，物体从 $x_1 = 0$ 到 $x_2 = \lambda_{\text{max}}$。
重力做功：$A_1 = mg\lambda_{\text{max}}$
弹性力做功：$A_2 = \int_{0}^{\lambda_{\text{max}}} (-kx)\,\mathrm{d}x = -\dfrac12 k\lambda_{\text{max}}^{2}$
始末动能均为零，由动能定理：
\[
mg\lambda_{\text{max}} - \frac12 k\lambda_{\text{max}}^{2} = 0 - 0
\]
\[
\lambda_{\text{max}} = \frac{2mg}{k}
\]
若缓慢放下（静平衡），$\lambda_{\text{sl}} = \dfrac{mg}{k}$，仅为 $\dfrac12\lambda_{\text{max}}$。

\item（书中例题 3.11，p.111，重点）长为 $l$ 的均质链条，部分置于水平面上，另一部分自然下垂，已知链条与水平面间静摩擦系数为 $\mu_0$，滑动摩擦系数为 $\mu$。
\begin{enumerate}
  \item 满足什么条件时，链条将开始滑动？
  \item 若下垂部分长度为 $b$ 时，链条自静止开始滑动，当链条末端刚刚滑离桌面时，其速度等于多少？
\end{enumerate}

\textbf{解：}设链条线密度为 $\rho$，下垂部分长度 $y$。
(1) 临界条件：下垂部分重力 $=$ 水平部分最大静摩擦力
\[
\rho y g = \mu_0 \rho (l - y)g \;\Longrightarrow\; y = \frac{\mu_0}{1 + \mu_0}l
\]
当 $y > \dfrac{\mu_0}{1 + \mu_0}l$ 时链条开始滑动。

(2) 以整个链条为研究对象，从下垂长度 $b$ 到链条刚滑离桌面（下垂长度 $l$）。
重力做功：
\[
A = \int_{b}^{l} \rho y g\,\mathrm{d}y = \frac12\rho g(l^{2} - b^{2})
\]
摩擦力做功：
\[
A' = -\int_{b}^{l} \mu\rho (l - y)g\,\mathrm{d}y = -\frac12\mu\rho g(l - b)^{2}
\]
由动能定理：
\[
\frac12\rho g(l^{2} - b^{2}) - \frac12\mu\rho g(l - b)^{2} = \frac12\rho l v^{2} - 0
\]
\[
v = \sqrt{\frac{g}{l}(l^{2} - b^{2}) - \frac{\mu g}{l}(l - b)^{2}}
\]

\item（书中例题 3.12，p.112）水平面内有一半径为 $R$ 的圆，在圆内离圆心 $O$ 距离为 $S$ 处有一质量很大、可视为固定的力心 $O'$，力心对单位质量的有心引力为 $\mu r$，$r$ 为力心到质量为 $m$ 的质点 $Q$ 的位矢大小，质点 $Q$ 被限制在圆周上运动。
\begin{enumerate}
  \item 求质点 $Q$ 从 $B$ 点由静止出发转过 $\phi$ 角有心力所做的功；
  \item 求质点通过第二象限所经历的时间。
\end{enumerate}

\textbf{解：}(1) 有心力 $F = m\mu r$。元功 $\mathrm{d}A = \vec{F}\cdot\mathrm{d}\vec{r} = FR\,\mathrm{d}\phi\sin\alpha$。
由正弦定理 $\dfrac{\sin\alpha}{S} = \dfrac{\sin\phi}{r}$，得 $\sin\alpha = \dfrac{S}{r}\sin\phi$。
\[
\mathrm{d}A = m\mu r \cdot R\,\mathrm{d}\phi \cdot \frac{S}{r}\sin\phi = m\mu RS\sin\phi\,\mathrm{d}\phi
\]
\[
A = \int_{0}^{\phi} m\mu RS\sin\phi\,\mathrm{d}\phi = m\mu RS(1 - \cos\phi)
\]

(2) 由动能定理：$A = \dfrac12 mR^{2}\left(\dfrac{\mathrm{d}\phi}{\mathrm{d}t}\right)^{2}$，代入 $A$ 得
\[
\frac12 mR^{2}\dot\phi^{2} = m\mu RS(1 - \cos\phi)
\]
\[
\dot\phi = \sqrt{\frac{2\mu S}{R}(1 - \cos\phi)} = 2\sqrt{\frac{\mu S}{R}}\sin\frac{\phi}{2}
\]
\[
\mathrm{d}t = \frac{\mathrm{d}\phi}{2\sqrt{\dfrac{\mu S}{R}}\sin\dfrac{\phi}{2}}
\]
通过第二象限：$\phi$ 从 $\dfrac{\pi}{2}$ 到 $\pi$。
\[
t = \int_{\pi/2}^{\pi} \frac{\mathrm{d}\phi}{2\sqrt{\dfrac{\mu S}{R}}\sin\dfrac{\phi}{2}}
= \frac{1}{2}\sqrt{\frac{R}{\mu S}}\int_{\pi/2}^{\pi} \frac{\mathrm{d}\phi}{\sin\dfrac{\phi}{2}}
\]
令 $u = \dfrac{\phi}{2}$，则
\[
t = \sqrt{\frac{R}{\mu S}}\int_{\pi/4}^{\pi/2} \frac{\mathrm{d}u}{\sin u}
= \sqrt{\frac{R}{\mu S}}\ln\tan\frac{u}{2}\Big|_{\pi/4}^{\pi/2}
= \sqrt{\frac{R}{\mu S}}\ln\tan\frac{\pi}{8}
\]
数值上 $t \approx 0.88\sqrt{\dfrac{R}{\mu S}}$。

\item 判断力 $\vec{F} = x^{2}y^{2}\,\vec{i} + x^{2}y^{2}\,\vec{j}$ 是不是保守力？$\vec{F} = 2xy\,\vec{i} + x^{2}\,\vec{j}$ 呢？

\textbf{解：}保守力判据：$\dfrac{\partial F_x}{\partial y} = \dfrac{\partial F_y}{\partial x}$。

(1) 对于 $\vec{F} = x^{2}y^{2}\,\vec{i} + x^{2}y^{2}\,\vec{j}$：
\[
\frac{\partial F_x}{\partial y} = 2x^{2}y,\qquad
\frac{\partial F_y}{\partial x} = 2xy^{2}
\]
两者不相等，故不是保守力。

(2) 对于 $\vec{F} = 2xy\,\vec{i} + x^{2}\,\vec{j}$：
\[
\frac{\partial F_x}{\partial y} = 2x,\qquad
\frac{\partial F_y}{\partial x} = 2x
\]
两者相等，故是保守力。

\item 把一个物体从地球表面上沿铅垂方向以第二宇宙速度 $v_0 = \sqrt{\dfrac{2GM}{R_e}}$ 发射出去，阻力忽略不计。求物体从地面飞行到与地心相距 $nR_e$ 处所经历的时间。

\textbf{解：}由机械能守恒：
\[
\frac12 mv_0^{2} - G\frac{M_e m}{R_e} = \frac12 mv^{2} - G\frac{M_e m}{x}
\]
代入 $v_0^{2} = \dfrac{2GM_e}{R_e}$，化简得
\[
v^{2} = \frac{2GM_e}{x}
\]
故 $v = \sqrt{\dfrac{2GM_e}{x}} = \dfrac{\mathrm{d}x}{\mathrm{d}t}$。
\[
\mathrm{d}t = \frac{\mathrm{d}x}{\sqrt{2GM_e}}\sqrt{x}
\]
积分：
\[
t = \int_{R_e}^{nR_e} \frac{\sqrt{x}}{\sqrt{2GM_e}}\,\mathrm{d}x
= \frac{1}{\sqrt{2GM_e}}\cdot\frac23 x^{3/2}\Big|_{R_e}^{nR_e}
= \frac{2}{3\sqrt{2GM_e}} R_e^{3/2}(n^{3/2} - 1)
\]

\item 用弹簧连接两个木板 $m_1$、$m_2$，弹簧压缩 $x_0$。求给 $m_2$ 上加多大的压力能使 $m_1$ 离开桌面？

\textbf{解：}取弹簧原长为弹性势能零点，地面为重力势能零点。初始压缩 $x_0 = \dfrac{m_2 g}{k}$。设加压力 $F$ 后弹簧再被压缩 $x_1 = \dfrac{F}{k}$，释放后弹簧伸长到 $x_2 = \dfrac{m_1 g}{k}$ 时 $m_1$ 恰好离开桌面。
机械能守恒：
\[
\frac12 k(x_0 + x_1)^{2} = \frac12 k x_2^{2} + m_2 g(x_0 + x_1 + x_2)
\]
代入 $x_0$、$x_1$、$x_2$ 的表达式，解得
\[
F = (m_1 + m_2)g
\]

\item（书中例题 3.15，p.126）物体 $M$ 悬于弹簧上，弹簧的弹性系数为 $k$，弹簧的原长与圆环的半径相等。不计摩擦力。求物体自弹簧的原长无初速度地沿圆环滑至最低点 $B$ 时所获得的动能。

\textbf{解：}取 $B$ 点为重力势能零点。弹簧原长与圆环半径 $R$ 相等。
初态（$A$ 点，弹簧原长）：
重力势能 $= mg(R + R\cos 60^{\circ}) = mg\cdot\frac32 R$，弹性势能 $= 0$，动能 $= 0$。
末态（$B$ 点）：
重力势能 $= 0$，弹性势能 $= \dfrac12 kR^{2}$，动能 $= E_k$。
机械能守恒：
\[
E_k + \frac12 kR^{2} = \frac32 mgR
\]
\[
E_k = \frac32 mgR - \frac12 kR^{2}
\]

\end{enumerate}

\end{document}
